\section{Kategorie 3: Governance (GOV)}
%==============================================================================

\epigraph{\textit{``The only freedom which deserves the name is that of pursuing our own good in our own way.''}}{--- John Stuart Mill, \textit{On Liberty} (1859)}

Die Governance-Kategorie (25 Axiome: MA-GOV-01 bis MA-GOV-25) definiert, \textbf{was erlaubt ist} -- sowohl im Sinne von Verhaltens-Constraints als auch im Sinne ethischer Prinzipien für Interventionen.

\subsection{Philosophischer Hintergrund: Von Mill zu Thaler}

\epigraph{\textit{``Libertarian paternalism is not an oxymoron.''}}{--- Thaler \& Sunstein (2003)}

Die Frage, unter welchen Bedingungen Verhaltensinterventionen ethisch legitim sind, berührt fundamentale Fragen der politischen Philosophie. Das Governance-Modul des \META{} navigiert zwischen zwei Extrempositionen: dem radikalen Libertarismus (``keine Intervention ist legitim'') und dem harten Paternalismus (``der Staat weiß es besser'').

\paragraph{Mills Harm Principle.}
John Stuart Mills \textit{On Liberty} (1859) etablierte das ``Harm Principle'' als liberales Grundprinzip: Die einzige Rechtfertigung für Einschränkung individueller Freiheit ist die Verhinderung von Schaden an anderen. Selbstschädigendes Verhalten darf nicht reguliert werden. Diese Position ist die historische Grundlage für die Ablehnung von Paternalismus.

\paragraph{Die paternalistische Herausforderung.}
Gegen Mill wurde argumentiert, dass Menschen systematisch gegen ihre eigenen Interessen handeln -- nicht aus freier Wahl, sondern aus kognitiven Limitationen. Gerald Dworkin (1972) unterschied:
\begin{itemize}
    \item \textbf{Harter Paternalismus}: Eingriff gegen den expliziten Willen der Person (z.B. Gurtpflicht)
    \item \textbf{Weicher Paternalismus}: Eingriff nur bei nicht-autonomen Entscheidungen (z.B. bei Manipulation, Zwang, oder fehlender Information)
\end{itemize}

\paragraph{Libertärer Paternalismus.}
Thaler \& Sunstein (2003, 2008) versöhnten die Positionen mit dem Konzept des ``Libertarian Paternalism'':
\begin{itemize}
    \item \textbf{Paternalistisch}: Interventionen zielen darauf, Menschen zu helfen, bessere Entscheidungen zu treffen (nach ihren eigenen Maßstäben)
    \item \textbf{Libertär}: Die Wahlfreiheit bleibt erhalten; es werden keine Optionen entfernt
\end{itemize}

Das Schlüsselargument: Es gibt \textit{keine neutrale} Choice Architecture. Jede Präsentation von Optionen beeinflusst die Wahl. Wenn Beeinflussung unvermeidbar ist, sollte sie wenigstens in eine ``gute'' Richtung gehen.

\paragraph{Die Kritik an Nudging.}
Bovens (2009) und andere haben Nudging kritisiert:
\begin{itemize}
    \item \textbf{Manipulation}: Nudges umgehen bewusste Deliberation
    \item \textbf{Transparenz}: Viele Nudges funktionieren nur, wenn sie nicht bekannt sind
    \item \textbf{Wer definiert ``gut''?}: Das Konzept setzt voraus, dass jemand weiß, was für andere gut ist
\end{itemize}

\paragraph{Die BCM-Position.}
Das \BCM{} ist \textbf{deskriptiv, nicht normativ}. Es modelliert Verhalten, es schreibt keine Ziele vor. Die Governance-Axiome definieren \textit{Constraints} für ethisches Interventionsdesign, aber die \textit{Ziele} werden extern vorgegeben. Diese Trennung von ``Ist'' und ``Soll'' folgt dem Humeschen Guillotine-Prinzip.

\subsection{Fundamentale Governance-Prinzipien}

\begin{axiom}[Autonomie-Primat: MA-GOV-03]
Individuelle Autonomie hat grundsätzlich Vorrang. Einschränkungen bedürfen der Rechtfertigung.

\paragraph{Philosophische Begründung.}
Autonomie -- die Fähigkeit zur Selbstgesetzgebung -- ist ein Kernwert der Aufklärung. Kant formulierte: ``Handle nur nach derjenigen Maxime, durch die du zugleich wollen kannst, dass sie ein allgemeines Gesetz werde.'' Autonome Entscheidungen verdienen Respekt, auch wenn sie ``schlecht'' erscheinen.

\paragraph{Operationalisierung.}
Im BCM bedeutet dies: Jede Intervention, die Autonomie einschränkt, muss \textit{explizit} gerechtfertigt werden. Die Beweislast liegt beim Intervenierenden.
\end{axiom}

\begin{axiom}[Nicht-Normativität: MA-GOV-04]
Das \BCM{} \textbf{beschreibt} Verhalten, es \textbf{schreibt nicht vor}. Die Wahl der Ziele liegt beim Auftraggeber, nicht beim Modell.

\paragraph{Philosophische Begründung.}
Dies folgt Humes Guillotine: Aus deskriptiven Aussagen (``so ist es'') folgen keine normativen (``so soll es sein''). Das BCM kann analysieren, \textit{wie} ein Ziel erreicht werden kann, aber nicht, \textit{ob} das Ziel gut ist.

\paragraph{Praktische Implikation.}
\BEATRIX{} kann sowohl für pro-soziale als auch für kommerzielle Ziele eingesetzt werden. Die ethische Verantwortung liegt beim Auftraggeber, nicht beim Modell.
\end{axiom}

\begin{axiom}[Transparenz-Gebot: MA-GOV-05]
Interventionen sollen transparent sein. Verdeckte Manipulation ist zu vermeiden.

\paragraph{Philosophische Begründung.}
Kants Publizitätsprinzip: ``Alle auf das Recht anderer Menschen bezogenen Handlungen, deren Maxime sich nicht mit der Publizität verträgt, sind unrecht.'' Eine Intervention, die nur funktioniert, wenn sie geheim bleibt, ist ethisch problematisch.

\paragraph{Empirisches Gegenargument.}
Einige Nudges funktionieren auch nach Offenlegung (Loewenstein et al., 2015). Aber nicht alle -- und hier entsteht ein ethisches Dilemma.
\end{axiom}

\begin{axiom}[Subsidiarität: MA-GOV-06]
Die \textbf{mildeste wirksame Intervention} ist zu wählen (Principle of Minimal Intervention).

\paragraph{Philosophische Begründung.}
Dies ist die Anwendung des Proportionalitätsprinzips auf Verhaltensinterventionen. Eine stärkere Intervention (z.B. Verbot) ist nur gerechtfertigt, wenn mildere Alternativen (z.B. Information) versagen.

\paragraph{Praktische Anwendung.}
Die 8 Interventionstypen des BCM sind nach ``Milde'' geordnet. Interventionsdesign sollte diese Hierarchie respektieren.
\end{axiom}

\subsection{Interventions-Taxonomie}

\begin{axiom}[8 Interventionstypen: MA-GOV-08]
Das \BCM{} unterscheidet acht fundamentale Interventionstypen:
\begin{enumerate}
    \item \textbf{NUDGE}: Änderung der Choice Architecture
    \item \textbf{BOOST}: Stärkung kognitiver Kompetenzen
    \item \textbf{THINK}: Aktivierung von System 2
    \item \textbf{INCENTIVE}: Materielle Anreize
    \item \textbf{REGULATE}: Ge- und Verbote
    \item \textbf{DESIGN}: Änderung des physischen Umfelds
    \item \textbf{SOCIAL}: Soziale Normen und Peer Effects
    \item \textbf{IDENTITY}: Aktivierung oder Änderung von Identitäten
\end{enumerate}
\end{axiom}

\begin{axiom}[Freiheits-Hierarchie: MA-GOV-09]
Die 8 Interventionstypen sind nach Freiheitsgrad geordnet:
\begin{equation}
\text{NUDGE} > \text{BOOST} > \text{THINK} > \text{INCENTIVE} > \text{SOCIAL} > \text{IDENTITY} > \text{DESIGN} > \text{REGULATE}
\end{equation}
wobei $>$ bedeutet ``mehr Wahlfreiheit erhaltend als''.
\end{axiom}

\begin{axiom}[Eskalations-Prinzip: MA-GOV-10]
Interventionen sollen von mild nach stark eskalieren:
\begin{equation}
\text{NUDGE} \rightarrow \text{BOOST} \rightarrow \text{INCENTIVE} \rightarrow \text{REGULATE}
\end{equation}
Erst wenn mildere Interventionen versagen, sind stärkere gerechtfertigt.
\end{axiom}

\subsection{Ethische Constraints}

\begin{axiom}[Distributive Gerechtigkeit: MA-GOV-12]
Kosten und Nutzen einer Intervention sind fair zu verteilen. Interventionen dürfen keine systematische Benachteiligung erzeugen.
\end{axiom}

\begin{axiom}[Reversibilität: MA-GOV-13]
Reversible Interventionen sind irreversiblen vorzuziehen.
\end{axiom}

\begin{axiom}[Crowding-Out vermeiden: MA-GOV-15]
\begin{equation}
\frac{\partial \text{Intrinsic}}{\partial \text{Extrinsic}} < 0 \quad \text{ist möglich}
\end{equation}
Extrinsische Anreize können intrinsische Motivation zerstören. Dies ist bei Interventionsdesign zu berücksichtigen.
\end{axiom}

%==============================================================================
